\section{Discusión}
La diferencia de conservación a favor de UMAP es más clara cuando se concede mayor peso a las escalas pequeñas. El objetivo de optimización local de UMAP ofrece una explicación compatible con este patrón. El diseño compara ajustes ya obtenidos y no aísla un mecanismo causal. El cambio de orden en vecindades más amplias impide convertir la AUC local en una afirmación sobre toda la geometría original.

La comparación PCA--ISOMAP muestra por qué importa distinguir el referente de cada métrica. La retención usa las identidades de los vecinos originales. La concordancia proyectada cambia ese referente y Procrustes compara la configuración completa de coordenadas tras alinearlas. La cercanía en la primera métrica y las diferencias en las otras son compatibles. Los caminos del grafo de ISOMAP alteran las distancias respecto a las euclidianas. Para interpretar esas alteraciones como recuperación de una variedad se necesitaría evidencia adicional sobre muestreo y estructura intrínseca.

La retención y la pureza son lógicamente distintas. La primera cuenta identidades conservadas y la segunda coincidencias de synset. En principio, cambiar vecinos por otros del mismo sentido puede preservar la pureza mientras se reduce la retención. En este piloto, UMAP tuvo la mayor AUC de conservación y el cambio de pureza ajustada puntualmente más favorable en la configuración principal. La evidencia simultánea respalda los dos contrastes de conservación frente a UMAP. Las bandas de los contrastes semánticos incluyen cero y dejan incierto su orden. El cambio de signo de $\Delta P_{\mathrm{aj}}$ al usar \texttt{abs5} muestra dependencia de la configuración. No establece un intercambio empírico entre conservar vecinos y preservar el sentido.

Las bandas muestran la estabilidad de la comparación frente a la reponderación documental con geometría fija. Que una banda incluya cero deja abiertas diferencias compatibles con sus límites. Evaluar equivalencia exigiría un margen especificado y un análisis que lo evaluara directamente.
